% ============================================================
%  SEKCJA 3: BIOLOGIA DOJRZEWANIA MIODU
% ============================================================

\section{Biologia dojrzewania miodu: proces ciągły, nie zdarzenie binarne}

Dojrzewanie miodu jest złożonym procesem biotransformacji, w~którym
nektar lub spadź ulega serii przemian enzymatycznych,
fizykochemicznych i~mikrobiologicznych, zanim osiągnie stan stabilności
jakościowej. Powszechne w~praktyce pszczelarskiej i~kontroli weterynaryjnej
utożsamianie dojrzałości miodu z~faktem zasklepienia komórek plastra woskiem stanowi uproszczenie nieadekwatne z~perspektywy biochemicznej. Zasklepienie jest bowiem jednym z~końcowych, morfologicznych
wskaźników dojrzewania, a~nie równoznaczne z~osiągnięciem pełnej
stabilności chemicznej i~mikrobiologicznej~\cite{zhang2021,sun2021}.

\subsection{Enzymatyczna transformacja nektaru}

Proces powstawania miodu inicjuje się już w~chwili pobierania nektaru lub spadzi
przez zbieraczki. Nektar, zdominowany przez sacharozę
i~zawierający wysoką zawartość wody, ulega pierwszej enzymatycznej
transformacji w~wolu miodowym (\textit{honey stomach crop}) owada,
gdzie miesza się z~wydzielinami gruczołów gardzielowych
pszczoły~\cite{biochemreactions2022}. Kluczowe enzymy inicjujące
transformację to:

\begin{itemize}
    \item \textbf{Inwertaza} ($\beta$-D-fruktozydaza), katalizuje
          hydrolizę sacharozy do~fruktozy i~glukozy; produkowana przez
          gruczoły podgarynkowe robotnic~\cite{biochemreactions2022};
          jej aktywność wzrasta stopniowo od~niskich wartości
          w~miodzie niedojrzałym do~wartości maksymalnych w~miodzie
          dojrzałym, co~czyni ją jednym z~najbardziej czułych
          wskaźników stopnia dojrzałości~\cite{zhang2021};
    \item \textbf{Oksydaza glukozowa}, utlenia glukozę do~kwasu
          glukonowego z~wydzieleniem nadtlenku wodoru (H$_2$O$_2$);
          kwas glukonowy odpowiada za~obniżenie pH miodu do~zakresu
          3{,}5--4{,}5 i~stanowi dominujący kwas organiczny
          w~dojrzałym miodzie; H$_2$O$_2$ pełni funkcję naturalnego
          czynnika antybakteryjnego~\cite{biochemreactions2022,zhang2021};
    \item \textbf{Diastaza} ($\alpha$- i~$\beta$-amylaza), katalizuje
          hydrolizę skrobi i~poliosacharydów nektarowych do~glukozy
          i~maltozy; aktywność diastazy jest parametrem normowanym
          przez Codex Alimentarius ($\geq 8$~j.~Schade) i~służy jako
          pośredni wskaźnik integralności termicznej i~dojrzałości
          miodu~\cite{codex_stan12,biochemreactions2022,zhang2021};
    \item \textbf{Katalaza}, rozkłada nadmiar H$_2$O$_2$, regulując
          równowagę między aktywnością antybakteryjną
          a~ochroną pozostałych składników miodu przed
          utlenianiem~\cite{biochemreactions2022}.
\end{itemize}

Badania aktywności enzymatycznej na~pięciu kolejnych etapach
transformacj, od~nektaru kwiatowego, przez wole zbieraczki
i~robotnic domowych, po~niezasklepione i~zasklepione komórki ula,  wykazały stopniowy wzrost aktywności inwertazy, amylazy oraz oksydazy glukozowej na~każdym kolejnym etapie
przetwarzania~\cite{enzymesrole2023}. Dane te potwierdzają,
że transformacja enzymatyczna przebiega jako \emph{kontinuum}
wzdłuż całego łańcucha produkcyjnego, a~nie jako zdarzenie
skokowe zachodzące w~momencie zasklepienia
komórki~\cite{enzymesrole2023,zhang2021}.

Równolegle do~przemian enzymatycznych pszczoły domowe aktywnie
odwadniają miód poprzez wentylację ula: uderzenia skrzydłami tworzą
stały przepływ powietrza nad otwartymi komórkami plastra, co~przyspiesza
odparowanie wody i~obniża wilgotność surowca z~wartości
charakterystycznych dla~nektaru do~finalnych ok.~18,5\%
w~miodzie przed zasklepieniem~\cite{zhang2021}. Pełne dojrzewanie
w~warunkach naturalnych trwa typowo kilkanaście dni,
przy czym czas ten może być znacznie wydłużony w~warunkach
wysokiej wilgotności względnej powietrza~\cite{zhang2021,tadele2025}.


\subsection{Dojrzewanie jako proces: dowody metabolomiczne}

Teza o~binarnym charakterze dojrzałości miodu, zgodnie z~którą
miód jest albo niedojrzały (niezasklepiony), albo dojrzały
(zasklepiony), znajduje coraz słabsze oparcie w~aktualnych
badaniach metabolomicznych. Sun i~in.~\cite{sun2021} w~badaniu
miodu rzepakowego~(\textit{Brassica napus}) wykazali, że~wskaźniki
dojrzałości, w~tym zawartość wody, stosunek fruktozy do~glukozy,
aktywność inwertazy i~diastazy oraz zawartość kwasów organicznych, zmieniają się \emph{gradientowo} w~trakcie procesu dojrzewania i~osiągają wartości stabilne niezależnie od~momentu zasklepienia
komórek przez pszczoły.

Guo i~in.~\cite{guo2020} przeprowadzili kompleksową analizę porównawczą profilu chemicznego i~aktywności biologicznej miodów dojrzałych i~niedojrzałych z~zastosowaniem metabolomiki opartej na~HPLC/QTOF/MS. Autorzy zidentyfikowali znaczące różnice w~zawartości polifenoli, aktywności antyoksydacyjnej (DPPH, FRAP) oraz profilu aminokwasów między miodem pobranym w~stanie niedojrzałości a~miodem dojrzałym, jednak różnice te~były ściśle skorelowane z~parametrami chemicznymi (zawartość wody, aktywność enzymatyczna), a~nie
wyłącznie z~faktem zasklepienia~\cite{guo2020}. Co~istotne, część próbek miodu pobranego przed pełnym zasklepieniem wykazywała profil metabolomiczny zbliżony do~dojrzałego, o~ile spełniała kryterium
niskiej wilgotności, co~stanowi bezpośrednie potwierdzenie koncepcji
dojrzałości funkcjonalnej. 

Sun i~in.~\cite{sun2021} przy użyciu GC-MS i~LC-MS zidentyfikowali
kwas dekenodwukarboksylowy (\textit{decenedioic acid}), kwas
tłuszczowy pochodzenia pszczelego, jako potencjalny biomarker
molekularny różnicujący miód dojrzały od~niedojrzałego. Biomarker
ten, walidowany dla~miodów rzepakowego, akacjowego i~jujube, był
istotnie bogatszy w~miodzie dojrzałym ($p < 0{,}01$), potwierdzając
że~dojrzewanie to~proces stopniowego gromadzenia specyficznych
metabolitów, nie~zaś binarne przejście~\cite{sun2021}.

Warto podkreślić, że~ekspozycja miodu na~wysoką temperaturę,
nawet po~pełnym zasklepieniu, może spowodować degradację enzymów
i~de~facto cofnięcie parametrów jakościowych do~wartości
charakterystycznych dla~miodu niedojrzałego~\cite{biochemreactions2022}.
Miód zasklepiony, poddany nieodpowiedniej obróbce termicznej,
może wykazywać aktywność diastazy poniżej progu Codex ($< 8$~j.~Schade)
i~podwyższoną zawartość HMF, co~dowodzi, że~morfologiczne kryterium
zasklepienia nie~gwarantuje utrzymania dojrzałości funkcjonalnej
w~całym łańcuchu dostaw.

Potwierdza to również praktyka pszczelarska, która fakt zasklepienia komórek traktuje jedynie jako wskaźnik potencjalnej dojrzałości miodu. Z jednej strony zdarza się, że w pewnych warunkach pszczoły zasklepiają komórki plastra, w których miód nie jest jeszcze dojrzały (zawartość wody przekracza 20% i krótko po wirowaniu miód fermentuje). Z drugiej strony powszechną praktyką jest pozyskiwanie miodu z plastrów poszytych częściowo (zazwyczaj przyjmuje się wskaźnik ¾ komórek zasklepionych w ramce) pod warunkiem, że w komórkach nie zasklepionych miód jest dojrzały. Faktycznym wskaźnikiem świadczącym o dojrzałości miodu jest zawartość wody, którą pszczelarz określa bądź metodą refraktometryczną bądź innymi testami (np. wstrząsanie ramki), których wynik jest bezpośrednio zależny od zawartości wody. 

\subsection{Koncepcja dojrzałości funkcjonalnej}
Na~podstawie powyższego przeglądu proponujemy operacyjną definicję
\textbf{dojrzałości funkcjonalnej miodu} (\textit{functional honey
maturity}), rozumianej jako zintegrowany stan produktu charakteryzujący
się jednoczesnym spełnieniem trzech grup kryteriów:

\begin{enumerate}
    \item \textbf{Stabilność mikrobiologiczna:} brak aktywnej
          fermentacji; aktywność wody $a_w < 0{,}60$; wilgotność
          $\leq 20\%$; liczba drożdży osmofilnych poniżej progu
          wzrostu w~danych warunkach temperatury
          i~przechowywania~\cite{zamora2006}.

    \item \textbf{Stabilność chemiczna:} stosunek fruktozy do~glukozy
          (F/G) $\geq 0{,}9$; zawartość HMF $\leq 40$~mg/kg;
          aktywność diastazy $\geq 8$~j.~Schade; aktywność inwertazy
          $\geq 4$~j.; zawartość kwasu glukonowego i~profil kwasów
          organicznych mieszczące się w~zakresie referencyjnym
          dla~danego typu botanicznego~\cite{codex_stan12,zhang2021}.

    \item \textbf{Stabilność biologiczna:} aktywność antyoksydacyjna
          (DPPH, FRAP) na~poziomie właściwym dla~danego typu
          botanicznego i~geograficznego; zawartość polifenoli i~flawonoidów
          nieobniżona w~stosunku do~wartości referencyjnych;
          aktywność antybakteryjna (oceniana przez~produkcję
          H$_2$O$_2$ i~obecność defensyn) niezmieniona~\cite{guo2020}.
\end{enumerate}

Definicja ta~celowo oddziela kryterium morfologiczne (zasklepienie
komórek) od~kryterium jakościowego, uznając, że~miód może osiągnąć
dojrzałość funkcjonalną przed pełnym zasklepieniem plastra (np.~przy
wilgotności 18\% w~umiarkowanym klimacie). Koncepcja ta~jest spójna
z~podejściem regulacyjnym Codex, który, jak wykazano w~sekcji~2 , penalizuje stan fermentacji, nie~morfologię~\cite{codex_stan12,sun2021}.
